\chapter{Une nouvelle famille}

Le lendemain matin, lorsque Khaer Ghal ouvrit les yeux, il lui fallut quelques
secondes pour se souvenir de l'endroit où il se trouvait.

Autour de lui, le campement s'éveillait déjà. Des voix se répondaient entre les
chariots, des chevaux remuaient dans leurs attaches et plusieurs voyageurs
rangeaient les affaires sorties la veille. Kamatsu dormait encore à côté de
lui, roulé dans la couverture qu'on leur avait donnée.

Khaer Ghal se redressa. La veille, personne ne leur avait demandé de partir.
On leur avait donné à manger, de l'eau et un endroit où dormir, mais rien de
plus ne leur avait été promis. Il ne savait donc pas encore ce qui allait se
passer.

Lorsque Kamatsu se réveilla à son tour, les deux garçons mangèrent ce qu'on
leur donna et restèrent quelque temps près du chariot où ils avaient passé la
nuit. Personne ne leur demanda de partir, mais personne ne leur expliqua non
plus ce qu'ils devaient faire.

Ils attendirent.

Puis la caravane reprit ses préparatifs. Les chevaux furent attelés, les
dernières caisses chargées et les feux éteints. Les voyageurs connaissaient
manifestement chacun leur rôle et circulaient entre les chariots avec
l'habitude de ceux qui avaient répété ces gestes des centaines de fois. Khaer
Ghal observait tout cela avec attention.

Lorsque les premiers chariots commencèrent finalement à avancer, un homme leur
fit simplement signe de suivre.

Ce ne fut pas une invitation solennelle. Personne ne leur annonça qu'ils
faisaient désormais partie de la caravane. On les laissa simplement marcher
avec elle.

Pour l'instant, cela suffisait.

\scenebreak

Les premiers jours furent étranges.

Kamatsu s'adapta rapidement. Il parlait facilement avec les voyageurs,
répondait à leurs questions et en posait probablement encore davantage. Très
vite, il connut plusieurs personnes par leur nom et semblait capable de
commencer une conversation avec presque n'importe qui.

Pour Khaer Ghal, les choses étaient différentes. Certains caravaniers
continuaient à le surveiller lorsqu'il passait près d'eux. Les regards
n'étaient plus aussi hostiles que lors de leur arrivée, mais la méfiance
n'avait pas complètement disparu. Son épée, ses défenses, sa peau gris-bleu et
son oreille abîmée suffisaient encore à rappeler à chacun d'où il venait.

Khaer Ghal ne leur en voulait pas vraiment. Lui aussi les observait.

Il apprenait leurs habitudes, regardait qui conduisait les chariots, qui
s'occupait des chevaux, qui montait la garde et qui préparait les repas. Il
découvrait peu à peu le fonctionnement de cette étrange communauté qui pouvait
changer de composition au fil des routes tout en continuant d'avancer comme un
seul groupe.

Au début, personne ne lui demandait grand-chose.

Puis les choses commencèrent à changer.

Un soir, alors que la caravane s'était arrêtée pour la nuit, Khaer Ghal vit un
homme essayer de déplacer seul une caisse trop lourde. Il s'approcha et attrapa
l'autre côté sans qu'on le lui demande.

L'homme le regarda avec surprise. Khaer Ghal attendit.

Après une seconde d'hésitation, le caravanier resserra simplement sa prise.

— Là-bas.

Ils transportèrent la caisse jusqu'à un chariot.

— Merci.

Le lendemain, il aida quelqu'un d'autre à ramasser du bois. Quelques jours
plus tard, il tint un cheval pendant qu'un voyageur vérifiait son sabot. Il
commença aussi à signaler certaines traces qu'il reconnaissait près des
campements ou les plantes qu'il savait dangereuses.

Ce n'étaient que de petites choses, mais elles s'accumulaient. Les caravaniers
découvrirent progressivement que le jeune Orc connaissait davantage la forêt
qu'ils ne l'avaient imaginé. Il savait reconnaître certaines empreintes,
repérer les signes laissés par plusieurs animaux et créatures et avait déjà
appris à se servir d'une arme.

Khaer Ghal, de son côté, découvrait que ce qu'il avait appris dans son village
pouvait servir à autre chose qu'à chasser ou combattre.

\scenebreak

Un après-midi, la caravane s'était arrêtée pour effectuer quelques réparations
avant de reprendre la route.

Khaer Ghal traversait le campement lorsqu'un homme occupé auprès d'un cheval
l'appela.

— Hé, Kher... Gal !

Khaer Ghal s'arrêta.

Ce n'était pas exactement son nom, mais il avait déjà compris que les sons qui
le composaient semblaient poser quelques difficultés aux gens de la caravane.
Il savait surtout parfaitement que l'homme s'adressait à lui et ne jugea pas
utile de le reprendre.

Le caravanier tenait une longue sangle de cuir appartenant au harnachement de
l'animal.

— Viens me tenir ça.

Khaer Ghal regarda l'homme, puis la sangle qu'il lui tendait. Pendant une
seconde, il resta immobile. Quelques jours auparavant, cet homme aurait
peut-être hésité avant de lui tourner le dos. À présent, il lui demandait
simplement de venir l'aider, sans méfiance particulière et sans même sembler
réfléchir à son geste.

— Khergal ?

Cette fois, Khaer Ghal s'approcha.

— Oui.

Il prit la sangle. Le caravanier retourna immédiatement à ce qu'il faisait et
ajusta une pièce du harnachement, vérifia rapidement son travail puis récupéra
la sangle.

— C'est bon.

Khaer Ghal la lui rendit.

— Merci.

Puis l'homme passa simplement de l'autre côté du cheval pour continuer son
travail. Khaer Ghal resta là une seconde avant de repartir.

Il n'y avait rien eu de remarquable dans cet échange. Personne ne s'était
arrêté pour les regarder et personne n'avait félicité le caravanier d'avoir
fait confiance à un Orc. L'homme avait eu besoin d'une deuxième paire de mains,
Khaer Ghal passait à proximité et il l'avait appelé.

C'était tout.

Et c'était précisément ce qui rendait ce moment différent.

\scenebreak

Les semaines passèrent.

La méfiance ne disparut pas en un jour, mais elle finit par s'user au rythme
des étapes, des repas et des tâches partagées. Khaer Ghal et Kamatsu cessèrent
progressivement d'être les deux enfants que l'on avait autorisés à suivre la
caravane pour devenir simplement deux enfants qui voyageaient avec elle.

Ils avaient désormais leur place autour des feux et savaient dans quel chariot
trouver une couverture lorsqu'il faisait froid. On leur demandait parfois
d'aider à charger ou décharger du matériel, d'aller chercher du bois ou de
surveiller quelque chose pendant quelques minutes. Kamatsu connaissait déjà
les histoires de plusieurs voyageurs et Khaer Ghal commençait à savoir vers
qui se tourner lorsqu'il voulait comprendre un geste, un outil ou une habitude
qu'il ne connaissait pas.

Bartiméus, lui, semblait avoir décidé depuis longtemps que toute cette période
d'adaptation ne le concernait absolument pas. Le chat apparaissait et
disparaissait au gré de ses envies, dormait dans les chariots qui lui
convenaient et venait réclamer de la nourriture avec une assurance qui
laissait penser qu'à ses yeux, la caravane lui appartenait probablement depuis
toujours.

Il venait également de plus en plus souvent retrouver les deux garçons.

Un soir, Kamatsu le regarda s'installer contre Khaer Ghal près du feu.

— Je crois qu'il t'aime bien.

Khaer Ghal baissa les yeux vers le chat.

— Il veut manger.

Bartiméus leva immédiatement la tête.

Kamatsu sourit.

— Tu vois ?

— Ça prouve rien.

— Il est venu vers toi le premier jour.

— Il voulait peut-être déjà manger.

Kamatsu éclata de rire. Khaer Ghal lui donna un léger coup d'épaule.

Autour d'eux, les conversations continuaient. Quelqu'un entretenait le feu,
deux voyageurs discutaient près d'un chariot et un autre terminait de ranger
du matériel avant la nuit. Personne ne prêtait particulièrement attention aux
deux garçons.

Khaer Ghal regarda un instant le campement. Quelques semaines auparavant, il
aurait surveillé chacun de ces mouvements, cherché les armes et évalué les
personnes qui l'entouraient. Désormais, il reconnaissait les voix, les
habitudes et même certaines disputes qui revenaient régulièrement.

Il n'était plus simplement toléré parmi eux : il était là comme Kamatsu, comme
Bartiméus et comme tous les autres.

Il lui fallut encore du temps pour comprendre ce que cela signifiait vraiment.
La caravane n'avait pas remplacé ce qu'ils avaient perdu et personne n'avait
cherché à le faire. Elle leur avait simplement offert quelque chose dont aucun
des deux garçons n'avait réellement osé se demander s'ils le retrouveraient un
jour : un endroit où revenir le soir, des visages familiers autour d'un feu et
des gens qui remarqueraient leur absence s'ils ne revenaient pas.

Peu à peu, sans cérémonie et presque sans qu'ils s'en aperçoivent, la caravane
devint leur nouvelle famille.